\documentclass[12pt,letterpaper]{article}
\renewcommand{\rmdefault}{cmr}
\renewcommand{\sfdefault}{cmss}
\renewcommand{\ttdefault}{cmtt}
\usepackage[T1]{fontenc}
\usepackage[utf8]{inputenc}
\usepackage[english,spanish]{babel}
\usepackage[letterpaper]{geometry}
\geometry{verbose,nomarginpar,tmargin=3cm,bmargin=3cm,lmargin=3cm,rmargin=3cm}
\setlength{\parskip}{\medskipamount}
\setlength{\parindent}{0pt}
\usepackage{array}
\usepackage{float}
\usepackage{longtable}
\usepackage{tabularx} % Paquete para tablas ajustadas
\usepackage{setspace}
\doublespacing

% Paquete de colores para las etiquetas de estado
\usepackage{xcolor}
\definecolor{colorIncluido}{RGB}{29, 78, 216}       % Azul claro / brillante
\definecolor{colorRedactado}{RGB}{30, 58, 138}      % Azul oscuro
\definecolor{colorParcial}{RGB}{107, 114, 128}      % Gris medio
\definecolor{colorConstruccion}{RGB}{156, 163, 175} % Gris claro
\definecolor{colorTextoGris}{RGB}{107, 114, 128}   % Gris notas de figuras/tablas

% Configuración de Hyperref (Citas y enlaces en negro)
\usepackage[unicode=true,
 bookmarks=true,bookmarksnumbered=false,bookmarksopen=false,
 breaklinks=false,pdfborder={0 0 0},backref=section,
 colorlinks=true,
 citecolor=black, % Citas bibliográficas en negro
 linkcolor=black, % Enlaces internos en negro
 urlcolor=black]  % Direcciones web en negro
 {hyperref}
\hypersetup{pdftitle={Titulo del proyecto sin tildes},
 pdfauthor={Nombres del los autores sin tildes},
 pdfsubject={tema de la investigacion},
 pdfkeywords={palabras claves}}

\makeatletter

%% Los convertidores HTML no saben que es tabularnewline
\providecommand{\tabularnewline}{\\}

% No generar la salida para \date
\date{}
% Utilizar el punto decimal, en vez de la acostumbrada coma del español.
\decimalpoint

% Etiqueta de figura y tabla a la izquierda seguida de punto.
\usepackage[justification=justified, singlelinecheck=false, labelsep = period]{caption}

% Centrar títulos de secciones
\usepackage{sectsty}
\sectionfont{\centering}

% Empezar sección en pagina impar
\let\stdsection\section
\renewcommand\section{\cleardoublepage\stdsection}

% Renombrar las secciones de acuerdo con NTC1486
\addto{\captionsspanish}{\renewcommand*{\contentsname}{\centering CONTENIDO}}
\addto{\captionsenglish}{\renewcommand*{\contentsname}{\centering CONTENIDO}}
\addto{\captionsspanish}{\renewcommand*{\listfigurename}{\centering LISTA DE FIGURAS}}
\addto{\captionsenglish}{\renewcommand*{\listfigurename}{\centering LISTA DE FIGURAS}}
\addto{\captionsspanish}{\renewcommand*{\listtablename}{\centering LISTA DE TABLAS}}
\addto{\captionsenglish}{\renewcommand*{\listtablename}{\centering LISTA DE TABLAS}}
\addto{\captionsspanish}{\renewcommand*{\refname}{REFERENCIAS}}
\addto{\captionsenglish}{\renewcommand*{\refname}{REFERENCIAS}}
\addto{\captionsspanish}{ \renewcommand*{\tablename}{Tabla}}
\addto{\captionsenglish}{ \renewcommand*{\tablename}{Tabla}}
\addto{\captionsenglish}{ \renewcommand*{\figurename}{Figura}}

% Crear un nuevo tipo de sección que sea centrado y sin numeración.
\makeatletter
\newcommand\@csection{\@startsection {section}{1}{\z@}%
  {-3.5ex \@plus -1ex \@minus -.2ex}%
  {2.3ex \@plus.2ex}%
  {\normalfont\Large\bfseries\centering}}
\newcommand{\csection}[1]{%
  \@csection*{#1}%
}
\makeatother

% Centrar la palabra BIBLIOGRAFÍA, sin modificar la entrada en la TOC.
\makeatletter
    \renewenvironment{thebibliography}[1]{%
    \csection{\refname} \bigskip
    \list{\@biblabel{\@arabic\c@enumiv}}
            {\settowidth\labelwidth{\@biblabel{#1}}%
    \leftmargin\labelwidth
    \advance\leftmargin\labelsep
           \@openbib@code
    \usecounter{enumiv}%
    \let\p@enumiv\@empty
    \renewcommand\theenumiv{\@arabic\c@enumiv}}%
    \sloppy
    \clubpenalty4000\@clubpenalty \clubpenalty
    \widowpenalty4000%
    \sfcode`\.\@m}
            {\def\@noitemerr{\@latex@warning{Empty `thebibliography' environment}}%
    \endlist}
\makeatother

% Comando para generar la línea punteada con la etiqueta de estado a la derecha
\newcommand{\tocitem}[3][0pt]{%
  \noindent\hspace{#1}{#2}\ \dotfill\ {#3}\par\vspace{-2pt}%
}

% Marcador visual para secciones pendientes
\newcommand{\enconstruccion}{%
  \vspace{0.5em}%
  \noindent\textcolor{colorConstruccion}{\textit{[En construcción: esta sección se redactará en las próximas sesiones del curso.]}}\par\vspace{1em}%
}


\begin{document}

% ==========================================
% PORTADA
% ==========================================
\begin{titlepage}
    \centering
    \begin{singlespace}
        % 1. ENCABEZADO INSTITUCIONAL
        UNIVERSIDAD PEDAGÓGICA Y TECNOLÓGICA DE COLOMBIA\\
        FACULTAD DE INGENIERÍA\\
        ESCUELA DE INGENIERÍA DE SISTEMAS Y COMPUTACIÓN

        \vfill % Distribución vertical elástica

        % 2. TÍTULO DEL PROYECTO
        \pdfbookmark[1]{Título}{TITULO}
        {\bfseries DEPENDENCIA DE NEQUI, DAVIPLATA Y BANCOLOMBIA
        PARA ACTIVIDADES COTIDIANAS EN ESTUDIANTES
        UNIVERSITARIOS DE LA UPTC, SEDE TUNJA}

        \vfill

        % 3. AUTORES

        Paula Nicol González Maldonado\\
        Brayan Smith Hernández Barrera\\
        Andrés Leonardo Herrera Agudelo


        \vfill

        % 4. CURSO Y CIUDAD/FECHA
        Documento elaborado en el marco del curso Metodología de la Investigación\\[0.5em]
        Tunja, Colombia $\cdot$ 2026
    \end{singlespace}
\end{titlepage}

\singlespacing

% ==========================================
% CONTENIDO
% ==========================================
\begin{center}
    {\Large\bfseries CONTENIDO}
\end{center}
\vspace{1em}

\tocitem{\hyperlink{sec:lot}{Índice de tablas}}{}
\tocitem{\hyperlink{sec:lof}{Índice de figuras}}{}
\tocitem{\hyperlink{sec:planteamiento}{\textbf{1. Planteamiento del problema}}}{}
\tocitem[1.5em]{\hyperlink{sec:antecedentes}{1.1 Antecedentes}}{}
\tocitem[1.5em]{\hyperlink{sec:descripcion}{1.2 Descripción de la situación problemática}}{}
\tocitem[1.5em]{\hyperlink{sec:formulacion}{1.3 Formulación del problema}}{}
\tocitem[1.5em]{\hyperlink{sec:consecuencias}{1.4 Consecuencias de no investigar}}{}
\tocitem[1.5em]{\hyperlink{sec:delimitacion}{1.5 Delimitación}}{}
\tocitem{\hyperlink{sec:justificacion}{\textbf{2. Justificación}}}{}
\tocitem{\hyperlink{sec:objetivos}{\textbf{3. Objetivos}}}{}
\tocitem{\hyperlink{sec:alcance}{\textbf{4. Alcance de la investigación}}}{}
\tocitem{\hyperlink{sec:hipotesis}{\textbf{5. Formulación de hipótesis}}}{}
\tocitem{\hyperlink{sec:variables}{\textbf{6. Variables}}}{}
\tocitem{\hyperlink{sec:marco_referencial}{\textbf{7. Marco referencial}}}{}
\tocitem{\hyperlink{sec:diseno_metodologico}{\textbf{8. Diseño metodológico}}}{}
\tocitem{\hyperlink{sec:resultados}{\textbf{9. Resultados y discusión}}}{}
\tocitem{\hyperlink{sec:conclusiones}{\textbf{10. Conclusiones y recomendaciones}}}{}
\tocitem{\hyperlink{sec:referencias}{Referencias}}{}
\tocitem{\hyperlink{sec:anexos}{Anexos}}{}

\vspace{2.5em}
% ==========================================
% ÍNDICE DE TABLAS
% ==========================================
\begin{center}
    {\Large\bfseries ÍNDICE DE TABLAS}
\end{center}
\vspace{0.5em}

\tocitem{Tabla 1. Delimitación de la investigación}{}
\tocitem{Tabla 2. Criterios de exclusión del estudio}{}
\tocitem{Tabla 3. Síntesis de las hipótesis de la investigación}{}
\tocitem{Tabla 4. Matriz de operacionalización y clasificación de variables}{}

\vspace{1.5em}

% ==========================================
% ÍNDICE DE FIGURAS
% ==========================================
\begin{center}
    {\Large\bfseries ÍNDICE DE FIGURAS}
\end{center}
\vspace{0.5em}
\begin{quote}
    \centering\color{colorTextoGris}\textit{Aún no hay figuras incorporadas al cuerpo del documento. Se listan aquí a medida que cada fase incorpore figuras propias (diagramas, gráficos, esquemas, etc.).}
\end{quote}

\vspace{2.5em}
\onehalfspacing


% ==========================================
% CUERPO DE LA INVESTIGACIÓN
% ==========================================
\hypertarget{sec:planteamiento}{}%
\begin{center}
    {\Large\bfseries 1. PLANTEAMIENTO DEL PROBLEMA}
\end{center}
\vspace{1em}

\hypertarget{sec:antecedentes}{}%
\noindent{\large\bfseries 1.1 ANTECEDENTES}\par\vspace{0.5em}

En los últimos años, el sector financiero global y nacional ha experimentado una acelerada transformación digital impulsada por el desarrollo de la tecnología financiera (fintech) y la masificación de los dispositivos móviles \cite{bermeo2022transformacion}. En Colombia, plataformas de servicios financieros digitales y billeteras electrónicas como Nequi, DaviPlata y las aplicaciones móviles de la banca tradicional (destacando Bancolombia) se han consolidado como herramientas fundamentales para la inclusión financiera, transformando radicalmente la forma en que las personas gestionan sus recursos monetarios, realizan transferencias y ejecutan pagos cotidianos sin necesidad de recurrir al efectivo \cite{banca2024reporte}.

Estudios sobre adopción tecnológica en jóvenes universitarios demuestran que este grupo poblacional es uno de los más receptivos y dependientes de las soluciones de banca móvil \cite{lopez2023adopcion}. Las dinámicas académicas, caracterizadas por transacciones inmediatas de bajo monto —como el pago de fotocopias, alimentación en cafeterías, transporte urbano y transferencias informales entre pares—, encuentran en las aplicaciones financieras un medio ágil que optimiza el tiempo y elimina las barreras asociadas al manejo de dinero físico o la asistencia a sucursales bancarias \cite{restrepo2023habitos}.

Sin embargo, el funcionamiento de estas plataformas digitales en entornos universitarios no ocurre en un plano aislado, sino que está estrechamente condicionado por la estabilidad de su infraestructura tecnológica, la disponibilidad de los servidores y la calidad de la conectividad \cite{martinez2023retos}. A diferencia de los centros urbanos con alta densidad comercial e infraestructura bancaria diversificada, las comunidades universitarias integradas en ciudades intermedias como Tunja presentan patrones de uso altamente concentrados. En el campus de la Universidad Pedagógica y Tecnológica de Colombia (UPTC), Sede Tunja, la masa estudiantil recurre de manera simultánea a las mismas plataformas en franjas horarias específicas, lo que genera una alta vulnerabilidad ante eventuales caídas del servicio o fallas operativas en los sistemas de las entidades \cite{pulido2024caracteristicas}.

Investigaciones previas sobre la adopción de tecnologías financieras suelen enfocar sus análisis en indicadores macroeconómicos de bancarización o en la usabilidad técnica de las interfaces de usuario desde la perspectiva del desarrollador de software \cite{silva2022evaluacion}. No obstante, existe una escasa documentación científica y empírica a nivel local que evalúe cuantitativa y cualitativamente el grado de dependencia funcional que han desarrollado los estudiantes universitarios frente a aplicaciones específicas como Nequi, DaviPlata y Bancolombia, así como la magnitud del impacto o parálisis en las actividades cotidianas que ocasiona la interrupción de estos servicios dentro del campus.

Esta brecha de conocimiento entre la penetración tecnológica de las billeteras digitales y la capacidad de contingencia real de los estudiantes universitarios ante la indisponibilidad del servicio constituye el punto de partida del presente estudio en la UPTC, Sede Tunja.

\vspace{1em}

\hypertarget{sec:descripcion}{}%
\noindent{\large\bfseries 1.2 DESCRIPCIÓN DE LA SITUACIÓN PROBLEMÁTICA}\par
En la cotidianidad de la Universidad Pedagógica Y Tecnológica de Colombia (UPTC), Sede Central Tunja, el uso de billeteras digitales e intermediarios financieros móviles como Nequi, DaviPlata y la aplicación de Bancolombia se ha integrado de forma generalizada en las dinámicas de consumo y movilidad de la comunidad estudiantil. La baja tenencia de dinero en efectivo entre los jóvenes universitarios ha provocado que transacciones esenciales tales como la compra de alimentos en las cafeterías del campus, el pago de fotocopias e impresiones de trabajos académicos, el pago del transporte público o colectivo urbano, e incluso la cancelación de arriendos en las zonas residenciales aledañas dependan casi de manera absoluta de la disponibilidad y correcto funcionamiento de estos sistemas informáticos.

A pesar de la alta conveniencia y agilidad que estas plataformas ofrecen, su infraestructura tecnológica presenta problemas operativos manifestados en caídas frecuentes de plataforma, interrupciones no programadas en los servidores, fallas de procesamiento durante picos de tráfico y congelamiento de saldos o transacciones pendientes. Cuando alguna de estas aplicaciones deja de operar o experimenta latencias prolongadas, el estudiante queda expuesto a un escenario de parálisis financiera inmediata. En el contexto de la UPTC Sede Tunja, esta situación genera problemas concretos como la imposibilidad de acceder a servicios básicos de alimentación durante los recesos de clase, retrasos en la entrega de actividades académicas físicas por no poder pagar la impresión de documentos, o dificultades de desplazamiento al no contar con efectivo para abordar el transporte urbano.


\vspace{1em}

\hypertarget{sec:formulacion}{}%
\noindent{\large\bfseries 1.3 FORMULACIÓN DEL PROBLEMA}\par
A partir de la situación problemática descrita, la pregunta general de investigación que orienta el presente estudio es:

¿De qué manera la disponibilidad y el estado operativo de las aplicaciones financieras digitales Nequi, DaviPlata y Bancolombia condicionan  las actividades cotidianas (académicas, de alimentación y movilidad) e inciden en el grado de dependencia de los estudiantes universitarios de la UPTC, Sede Tunja?

Para abordar de manera sistemática la pregunta central, la investigación se desglosa en las siguientes sub-preguntas:

1.	¿Cuál es el nivel de dependencia de los estudiantes frente a cada una de estas aplicaciones y cuáles son las actividades cotidianas más afectadas cuando se presentan interrupciones en el servicio?

2.	¿Existe un contraste significativo en el grado de vitalidad e impacto operacional entre Nequi, DaviPlata y Bancolombia durante una falla o caída de sistema?

3.	¿Qué mecanismos de contingencia (como uso de efectivo de reserva, migración a plataformas alternativas o cancelación de compras) implementan los estudiantes de la UPTC ante la indisponibilidad de su billetera digital preferida?


\vspace{1em}

\hypertarget{sec:consecuencias}{}%
\noindent{\large\bfseries 1.4 CONSECUENCIAS DE NO INVESTIGAR}\par
De no realizarse esta investigación, continuara la parálisis que sufren los estudiantes de la UPTC ante las caídas de Nequi, DaviPlata y Bancolombia, invisibilizando la alta vulnerabilidad a la que se ven expuestos en sus actividades diarias de alimentación, transporte y cumplimiento académico. Asimismo no permitirá el desarrollo de los comercios del campus y a las entidades bancarias de insumos fundamentados para promover alternativas de pago o exigir mayores estándares de disponibilidad en plataformas financieras digitales.

\vspace{1em}

\hypertarget{sec:delimitacion}{}%
\noindent{\large\bfseries 1.5 DELIMITACIÓN}\par

La presente investigación se enfoca en caracterizar el nivel de dependencia hacia las billeteras digitales (Nequi, DaviPlata y Bancolombia) y evaluar el impacto que generan sus caídas operativas en las actividades cotidianas de los estudiantes de la UPTC, Sede Tunja. En el Cuadro 1 se detallan las dimensiones de la delimitación aplicada al estudio.

\begin{table}[htbp]
\centering
\begin{tabular}{|p{3cm}|p{12cm}|}
\hline
\textbf{Dimensión} & \textbf{Delimitación aplicada} \\ \hline
Temática & Análisis de la dependencia funcional e impacto operacional causado por fallas en tres aplicaciones de banca móvil (Nequi, DaviPlata y Bancolombia) en actividades cotidianas (alimentación, transporte, pagos académicos); no abarca el desarrollo de software ni auditorías de seguridad informática. \\ \hline
Poblacional & Muestra de estudiantes universitarios de pregrado matriculados en la UPTC, Sede Tunja, que utilicen activamente al menos una de las tres plataformas analizadas. \\ \hline
Espacial & Campus central de la Universidad Pedagógica y Tecnológica de Colombia (UPTC), Sede Tunja, incluyendo comercios internos (cafeterías, centros de fotocopiado) y zonas residenciales aledañas. \\ \hline
Temporal & Marco temporal de 16 semanas, correspondiente al semestre lectivo del año 2026. \\ \hline
\end{tabular}
\caption{Delimitación de la investigación}
\label{tab:delimitacion}
\end{table}

\begin{table}[htbp]
\centering
\begin{tabular}{|p{4.5cm}|p{10.5cm}|}
\hline
\textbf{Queda fuera} & \textbf{Por qué} \\ \hline
Otras aplicaciones financieras o criptomonedas & Se evalúan únicamente Nequi, DaviPlata y Bancolombia por ser las de mayor penetración y uso dentro de la comunidad universitaria local. \\ \hline
Estudiantes de otras sedes (Sogamoso, Duitama, Chiquinquirá) & Las dinámicas de movilidad, acceso comercial y uso del campus varían según la sede; el estudio se concentra en Tunja. \\ \hline
Análisis del código de programación de las aplicaciones & El estudio aborda la perspectiva del usuario y el impacto operativo, no el desarrollo o mantenimiento del software bancario. \\ \hline
\end{tabular}
\caption{Criterios de exclusión del estudio}
\label{tab:exclusion}
\end{table}

\hypertarget{sec:justificacion}{}%
\begin{center}
    {\Large\bfseries 2. JUSTIFICACIÓN}
\end{center}
\vspace{1em}

Este estudio busca analizar y comparar de qué manera la disponibilidad y el estado operativo de Nequi, DaviPlata y Bancolombia condicionan las actividades cotidianas e inciden en el grado de dependencia de los estudiantes de la UPTC, Sede Tunja. Dentro de la problemática ya identificada sobre la baja tenencia de efectivo y la alta vulnerabilidad en el campus, esta investigación cumple el papel de pasar de la percepción informal a la evidencia empírica. Con ello, se espera obtener conocimiento cuantitativo inédito sobre la magnitud del impacto operativo de las fallas, la clasificación de criticidad de cada plataforma y el nivel real de preparación transaccional de los estudiantes, información que hoy no existe a nivel local.

Es necesario investigarlo ahora porque la adopción de billeteras digitales en Colombia se ha masificado aceleradamente y las caídas o fallas masivas de estas aplicaciones son un fenómeno real y recurrente en el país. En el campus de la UPTC Sede Tunja, los estudiantes realizan múltiples pagos de bajo monto cada semana (alimentación, fotocopias y transporte) bajo horarios fijos.

El estudio sirve para que los estudiantes reconozcan su grado de dependencia y la importancia de contar con alternativas de pago en su rutina académica. A su vez, para la Ingeniería de Sistemas representa un aporte relevante al evaluar las consecuencias de la indisponibilidad de software en la vida real, subrayando la necesidad de desarrollar sistemas de información transaccionales con alta disponibilidad, tolerancia a fallas y arquitecturas resilientes ante la alta demanda.

La investigación se aborda mediante un enfoque descriptivo y comparativo sustentado en un instrumento estandarizado de recolección de datos, el cual resulta más conveniente que una simple encuesta de percepción aislada porque permite relacionar directamente eventos de falla operacionales con consecuencias medibles en la cotidianidad de los usuarios. Además, este diseño metodológico deja una matriz de evaluación replicable que servirá de base para futuras investigaciones sobre usabilidad, disponibilidad y dependencia de servicios digitales en otras instituciones de educación superior de la región.

\hypertarget{sec:objetivos}{}%
\begin{center}
    {\Large\bfseries 3. OBJETIVOS}
\end{center}
\vspace{1em}

\noindent{\large\bfseries 3.1 OBJETIVO GENERAL}\par\vspace{0.5em}
Analizar de qué manera la disponibilidad y el estado operativo de las aplicaciones financieras digitales Nequi, DaviPlata y Bancolombia condicionan la ejecución de las actividades cotidianas e inciden en el grado de dependencia de los estudiantes universitarios de la UPTC, Sede Tunja, mediante la aplicación de un estudio, para identificar las principales afectaciones ante fallas del servicio y reconocer los mecanismos de contingencia implementados.

\vspace{1em}
\noindent{\large\bfseries 3.2 OBJETIVOS ESPECÍFICOS}\par\vspace{0.5em}
\begin{itemize}
    \item[\textbf{a)}] \textbf{Identificar} el nivel de dependencia de los estudiantes respecto a Nequi, DaviPlata y Bancolombia, determinando las actividades cotidianas específicas que resultan más afectadas ante una caída o falla del servicio.\\
    \textit{Meta:} Aplicación de un instrumento de recolección de datos y elaboración de una matriz de caracterización de frecuencia de uso, hábitos de pago y nivel de afectación operacional en los estudiantes de la UPTC, Sede Tunja.

    \item[\textbf{b)}] \textbf{Comparar} Comparar cuál de las tres plataformas analizadas (Nequi, DaviPlata y Bancolombia) presenta un mayor grado de vitalidad e impacto operacional en las actividades diarias de la comunidad estudiantil evaluada.\\
    \textit{Meta:} Matriz comparativa y reporte cuantitativo de vulnerabilidad que clasifique e identifique la aplicación de mayor criticidad para el desarrollo de la rutina universitaria.

    \item[\textbf{c)}] \textbf{Reconocer} los mecanismos de contingencia y las alternativas (uso de efectivo de reserva, migración a otras plataformas o cancelación de pagos) que implementan los estudiantes ante la indisponibilidad de su billetera digital preferida.\\
    \textit{Meta:} Fichas de caracterización de respuestas de contingencia y mapa de hábitos de respaldo frente a la falla de sistemas financieros móviles en el campus universitario.

\end{itemize}
\hypertarget{sec:alcance}{}%
\begin{center}
    {\Large\bfseries 4. ALCANCE DE LA INVESTIGACIÓN}
\end{center}
\vspace{1em}

\noindent Este estudio es de tipo \textbf{descriptivo, con un componente comparativo}: busca caracterizar, a partir de una encuesta aplicada a los estudiantes de pregrado de la UPTC, Sede Tunja, el nivel de dependencia funcional hacia Nequi, DaviPlata y Bancolombia, las actividades cotidianas (alimentación, transporte y pagos académicos) que resultan afectadas cuando alguna de ellas falla y los mecanismos de contingencia que los estudiantes emplean, además de contrastar en qué medida esa dependencia varía entre las tres aplicaciones. Se elige este alcance, y no uno exploratorio, porque el fenómeno ya está definido y documentado: en 2023, según Colombia Fintech, 13 billeteras digitales reunieron 54,63 millones de usuarios en el país (cifra que incluye duplicidades, pues una misma persona usa varias billeteras) y procesaron 3.342 millones de transacciones por cerca de 332 billones de pesos, con Nequi como líder del segmento, seguida de Bancolombia a la Mano y DaviPlata~\cite{colfintech2024}, mientras que la caída total de los servicios digitales de Bancolombia y Nequi del 24 de octubre de 2025~\cite{elcolombiano2025} y la falla masiva de Bancolombia del 22 y 23 de febrero de 2026, que generó cerca de 8.000 quejas ante la Superfinanciera en 48 horas~\cite{infobae2026}, evidencian que las caídas del servicio son reales y recurrentes; tampoco se plantea un alcance explicativo, porque, como señala el numeral 1.1, la documentación local sobre el grado de dependencia estudiantil es escasa y el diseño no manipula ni controla las caídas del servicio, por lo que describe y compara, pero no atribuye causalidad.
\par\vspace{1.2em}

\noindent\textbf{Limitaciones de la investigación}
\par\vspace{0.6em}

\noindent Al basarse en una muestra no probabilística de estudiantes de una sola sede, los resultados describen a los participantes y no pueden generalizarse estadísticamente a toda la comunidad de la UPTC, a otras sedes ni a otras universidades; los datos provienen de autorreportes, por lo que la dependencia percibida puede diferir del uso real y los relatos sobre caídas pasadas están expuestos a sesgo de recuerdo; las interrupciones del servicio no se provocan ni se observan de forma controlada, y variables externas como la duración y el tipo de falla, la hora del día, el saldo disponible o la calidad de la conectividad no pueden controlarse del todo, de modo que el estudio no permite establecer relaciones causales; dado que un mismo estudiante puede usar más de una de las tres aplicaciones, las comparaciones entre ellas no se hacen entre grupos completamente independientes, lo que debe considerarse al elegir las pruebas estadísticas; y el instrumento de recolección se valida dentro de este mismo estudio, por lo que sus propiedades (validez de contenido y confiabilidad) deben reportarse junto con los resultados finales, no darse por supuestas.
\par\vspace{0.8em}
\hypertarget{sec:hipotesis}{}%
\begin{center}
    {\Large\bfseries 5. FORMULACIÓN DE HIPÓTESIS}
\end{center}
\vspace{1em}

\noindent Dado que el alcance del estudio es descriptivo con un componente comparativo secundario (numeral 4), las hipótesis se formulan de dos tipos: \textit{descriptivas} (H1a y H1b), que anticipan la tendencia de una variable sin relacionarla con otra, y de \textit{diferencia de grupos} (H1c), que anticipa que la afectación difiere entre las aplicaciones. No se formulan hipótesis causales, pues el diseño no manipula ni controla las caídas del servicio. Cada hipótesis de trabajo (H1) se acompaña de su hipótesis nula (H0), que es la que se pone a prueba en el análisis de resultados: no se demuestra H1 de forma directa, sino que se evalúa si existe evidencia suficiente para rechazar H0. Los objetivos específicos a) y e), de caracterización y de recomendaciones, no requieren hipótesis propia.
\par\vspace{1.2em}

\noindent Las hipótesis no fijan un porcentaje, porque los antecedentes disponibles no permiten justificar un valor de corte; la proporción observada se estimará y se reportará en los resultados. Las hipótesis se refieren a los participantes de la muestra; cualquier generalización a una población más amplia deberá definirse en el diseño metodológico, junto con la población a la que se refiere y los supuestos que la sostienen. Se redactan de modo que un resultado sin diferencia, o en sentido contrario, también constituya un hallazgo válido.
\par\vspace{1.2em}

\noindent\textbf{H1a (hipótesis de trabajo --- nivel de uso):} En la muestra predominan los estudiantes de pregrado de la UPTC, Sede Tunja, que realizan la mayoría o casi todos los pagos de sus actividades cotidianas (alimentación, transporte, fotocopias y pagos académicos) con al menos una de las tres aplicaciones (Nequi, DaviPlata o Bancolombia).
\par\vspace{0.8em}

\noindent\textbf{H0a (hipótesis nula --- nivel de uso):} En la muestra no predominan los estudiantes que realizan la mayoría o casi todos los pagos de sus actividades cotidianas con al menos una de las tres aplicaciones.
\par\vspace{1.2em}

\noindent\textbf{H1b (hipótesis de trabajo --- consecuencias de la caída):} Entre los estudiantes de la muestra que vivieron en el semestre una caída o falla de alguna de las tres aplicaciones, confirmada como problema de la aplicación y no de su conexión o dispositivo, predominan quienes reportan al menos un tipo de afectación en alguna actividad cotidiana (aplazar o retrasar la actividad, cancelar una compra o un pago, no cumplir una actividad académica a tiempo o tener un costo adicional).
\par\vspace{0.8em}

\noindent\textbf{H0b (hipótesis nula --- consecuencias de la caída):} Entre los estudiantes de la muestra que vivieron una caída o falla confirmada, no predominan quienes reportan al menos un tipo de afectación en alguna actividad cotidiana.
\par\vspace{1.2em}

\noindent\textbf{H1c (hipótesis de trabajo --- comparación secundaria entre aplicaciones):} Entre los estudiantes de la muestra que usan cada aplicación, la afectación de las actividades financieras por su no disponibilidad difiere según la aplicación (Nequi, DaviPlata o Bancolombia).
\par\vspace{0.8em}

\noindent\textbf{H0c (hipótesis nula --- comparación secundaria entre aplicaciones):} Entre los estudiantes de la muestra que usan cada aplicación, la afectación de las actividades financieras por su no disponibilidad no difiere según la aplicación.
\par\vspace{1.2em}

\noindent La Tabla~\ref{tab:hipotesis-resumen} sintetiza la relación entre cada par de hipótesis, su tipo, el objetivo específico que responde y las variables que involucra.

\begin{table}[!ht]
\centering
\small
\renewcommand{\arraystretch}{1.3}
\begin{tabular}{|c|p{3.4cm}|p{3.6cm}|p{4cm}|}
\hline
\textbf{Hipótesis} & \textbf{Tipo} & \textbf{Objetivo} & \textbf{Variable(s) involucrada(s)} \\ \hline
H1a / H0a & Descriptiva & b) Nivel de uso & Uso \\ \hline
H1b / H0b & Descriptiva & c) Consecuencias de la no disponibilidad & Falla reportada y afectación \\ \hline
H1c / H0c & Diferencia de grupos (comparativa secundaria) & d) Cuantificar la afectación & Aplicación y afectación \\ \hline
\end{tabular}
\caption{Síntesis de las hipótesis de la investigación.}
\label{tab:hipotesis-resumen}
\end{table}

\hypertarget{sec:variables}{}%
\begin{center}
    {\Large\bfseries 6. VARIABLES}
\end{center}
\vspace{1em}

Las hipótesis del numeral 5 relacionan las siguientes variables, que se identifican, se clasifican según su función y se definen conceptualmente a partir del marco teórico. Su definición operacional (indicadores, escalas de medición, fuentes de datos y procedimiento de cálculo) se desarrollará en el diseño metodológico. Como el alcance es descriptivo con componente comparativo (numeral 4) y no se manipulan las caídas del servicio, la clasificación por función es analítica, no causal.

\begingroup
\small
\begin{longtable}{|p{3.5cm}|p{2.5cm}|p{6.5cm}|p{1.5cm}|}
\caption{Matriz de operacionalización y clasificación de variables} \label{tab:matriz_variables} \\
\hline
\textbf{Variable} & \textbf{Tipo} & \textbf{Definición Operacional / Conceptual} & \textbf{Hipótesis} \\ \hline
\endfirsthead

\multicolumn{4}{c}{{\thetable{} }} \\
\hline
\textbf{Variable} & \textbf{Tipo} & \textbf{Definición Operacional / Conceptual} & \textbf{Hipótesis} \\ \hline
\endhead



\hline
\endlastfoot

Nivel de dependencia funcional hacia la aplicación financiera digital & Dependiente (principal) & Grado en que el estudiante necesita una billetera digital o app bancaria móvil (Nequi, DaviPlata o Bancolombia) para realizar sus pagos cotidianos de alimentación, transporte y pagos académicos. Se manifiesta en la frecuencia de uso, la proporción de pagos que se hacen por este medio y la dificultad de reemplazarlo. & H1a, H1b \\ \hline
Aplicación financiera digital utilizada & Independiente (categórica) & Plataforma de pago móvil que el estudiante usa activamente: Nequi, DaviPlata o Bancolombia. Es el factor de comparación entre grupos. Un estudiante puede usar más de una, por lo que los grupos no son totalmente independientes. & H1b \\ \hline
Impacto operacional de la interrupción del servicio & Dependiente & Grado en que una caída, falla o latencia prolongada de la aplicación altera las actividades cotidianas del estudiante: aplazarlas, cancelarlas o modificarlas, y qué tan grave es la afectación percibida. & H1c \\ \hline
Mecanismos de contingencia ante la indisponibilidad & Dependiente & Estrategias con que el estudiante sigue con sus actividades cuando su billetera preferida no funciona: efectivo de reserva, cambio a otra plataforma, ayuda de pares o cancelación de la compra o el pago. & H1d \\ \hline
Tipo de actividad cotidiana afectada & Interviniente & Categoría de actividad que depende del pago digital: alimentación en cafeterías, transporte urbano, fotocopias e impresiones, arriendo u otras. La afectación varía según la urgencia y la posibilidad de sustituir el medio de pago. & H1a, H1c \\ \hline
Experiencia previa con caídas del servicio & Interviniente & Haber vivido o no una interrupción de la aplicación, y con qué frecuencia y duración. Condiciona la percepción del impacto y la preparación de planes de contingencia. Define además la subpoblación de H1c. & H1c, H1d \\ \hline
Hábitos de pago y disponibilidad de efectivo & Interviniente & Preferencia por pagos digitales frente al efectivo, y monto o hábito de mantener dinero físico. Modula tanto el nivel de dependencia como la capacidad de contingencia. & H1a, H1d \\ \hline
Características sociodemográficas y académicas & Interviniente / de control & Programa, semestre, edad, género y situación de vivienda (residente o no en zonas aledañas al campus), que pueden explicar diferencias en el uso y la dependencia. & H1a, H1b \\ \hline
\end{longtable}
\endgroup
\begin{center}
    {\Large\bfseries 7. MARCO REFERENCIAL}
\end{center}

\noindent\textbf{\large 7.1 Marco Histórico}\par\vspace{0.2cm}
La dependencia de los estudiantes universitarios hacia aplicaciones como Nequi, DaviPlata y Bancolombia es el resultado de una rápida transformación del ecosistema financiero colombiano que, en menos de una década, reemplazó el efectivo por los depósitos de bajo monto \cite{unisalle2026}. Este fenómeno histórico inició con la eliminación de las barreras bancarias tradicionales (papeleo, cuotas de manejo y sucursales físicas) a través de plataformas digitales, marcando un hito fundacional con el lanzamiento de Nequi en abril de 2016 y la expansión de DaviPlata \cite{unisalle2026}. Sin embargo, la adopción de estas tecnologías experimentó su aceleración definitiva con la pandemia de COVID-19, un evento que actuó como catalizador: el 61 \% de los colombianos disminuyó la frecuencia de uso del dinero físico y el volumen de pagos electrónicos aumentó un 70 \% ante las nuevas necesidades de bioseguridad \cite{cajamarca2021}.

Esta transición cualitativa provocó que, para el cierre de 2024, el 76,1 \% de la población adulta en Colombia ya utilizara depósitos de bajo monto, un producto que en poco tiempo logró que el 84,1 \% de sus cuentas se mantuvieran activas con un promedio de 92 operaciones anuales por usuario (casi el doble de la actividad registrada en las cuentas de ahorro tradicionales) \cite{unisalle2026}. La digitalización masiva de los micropagos mediante códigos QR y transferencias de número a número se afianzó aún más en la cotidianidad con la entrada en operación del sistema Bre-B en octubre de 2025 \cite{unisalle2026}. Es al final de esta evolución histórica donde la comunidad estudiantil de la UPTC, Sede Tunja, adaptó sus dinámicas de consumo —como el pago de fotocopias, alimentación en cafeterías y transporte urbano— a estas billeteras, llegando al punto de vulnerabilidad actual donde las caídas operativas de la infraestructura tecnológica paralizan su rutina académica y financiera \cite{restrepo2023habitos, martinez2023retos, pulido2024caracteristicas}.

\vspace{0.4cm}
\noindent\textbf{\large 7.2 Estado del Arte}\par\vspace{0.2cm}
\textit{Adopción y hábitos de uso de billeteras digitales en universitarios:} La literatura reciente coincide en la alta penetración de las plataformas financieras en la población joven. López y Vargas \cite{lopez2023adopcion} analizan la adopción desde el modelo UTAUT, destacando que los universitarios son altamente receptivos a la banca móvil por su facilidad de uso. En una línea similar, Restrepo y Gómez \cite{restrepo2023habitos} evalúan los hábitos de consumo en educación superior, coincidiendo en que las aplicaciones se utilizan predominantemente para transacciones inmediatas de bajo monto, como fotocopias y cafeterías. Ambos estudios concuerdan en la masificación de estas herramientas, pero se enfocan exclusivamente en los factores que motivan su adopción y uso en condiciones normales, sin contemplar escenarios de interrupción del servicio.

\textit{Infraestructura técnica e impacto operacional:} Desde una perspectiva más técnica, la literatura aborda las vulnerabilidades de estas plataformas. Martínez \cite{martinez2023retos} documenta los retos de disponibilidad en la infraestructura de la banca móvil colombiana, señalando la recurrencia de caídas en los servidores. Por su parte, Silva y Torres \cite{silva2022evaluacion} evalúan el impacto operacional de las aplicaciones de pago, concluyendo que las fallas de software reducen drásticamente la usabilidad percibida. Aunque ambos autores convergen en que la inestabilidad técnica es un problema latente, sus enfoques se centran en el rendimiento del software y en indicadores técnicos.

\textit{Identificación del vacío de investigación:} Al cruzar ambos grupos de estudios, se evidencia un vacío claro: la literatura actual documenta exhaustivamente la adopción tecnológica (desde el usuario) y los problemas de infraestructura (desde el software), pero ninguno de los estudios revisados aborda el cruce de ambas variables en un entorno microeconómico cautivo; es decir, no se evalúa el grado de parálisis o el impacto directo en las actividades cotidianas vitales (alimentación, movilidad, academia) de una comunidad estudiantil cuando ocurren estas caídas operativas. Ese es el vacío que aborda este proyecto, cuantificando empíricamente la dependencia funcional y la capacidad de contingencia en los estudiantes de la UPTC, Sede Tunja.

\vspace{0.4cm}
\noindent\textbf{\large 7.3 Marco Teórico}\par\vspace{0.2cm}
El marco teórico de esta investigación se sustenta en dos modelos conceptuales que permiten explicar tanto la adopción masiva de las plataformas financieras como el comportamiento de los estudiantes ante sus fallas operativas.

\begin{enumerate}
    \item \textit{Teoría Unificada de Aceptación y Uso de Tecnología (UTAUT):} Para abordar la primera y segunda sub-pregunta de investigación sobre el nivel de dependencia hacia Nequi, DaviPlata y Bancolombia, este estudio se apoya en el modelo UTAUT. Esta teoría explica por qué los usuarios adoptan un sistema de información basándose en la expectativa de rendimiento (la agilidad en los pagos), la expectativa de esfuerzo (facilidad de uso) y las condiciones facilitadoras (la infraestructura del campus). Bajo este modelo, la dependencia funcional que experimentan los estudiantes de la UPTC no es un fenómeno aleatorio, sino el resultado de que estas plataformas cumplen con alta eficiencia sus necesidades de transacciones cotidianas de bajo monto, generando un anclaje tecnológico profundo.

    \item \textit{Teoría de la Resiliencia de Sistemas y Dependencia del Medio:} Para abordar la tercera sub-pregunta, referente al impacto operacional y los mecanismos de contingencia frente a las caídas del servicio, se utiliza la intersección entre la Teoría de Sistemas (desde la resiliencia en la ingeniería) y la Teoría de la Dependencia del Medio. Esta perspectiva explica por qué una disrupción técnica en un nodo del sistema (el servidor del banco) genera un efecto en cascada que paraliza subsistemas físicos (alimentación, transporte urbano). Asimismo, permite interpretar, desde la elección racional, por qué los usuarios implementan (o fallan al implementar) mecanismos de contingencia, como el efectivo de reserva, buscando recuperar el equilibrio en sus actividades cotidianas ante la indisponibilidad de la infraestructura digital principal.
\end{enumerate}

\vspace{0.4cm}
\noindent\textbf{\large 7.4 Marco Conceptual}\par\vspace{0.2cm}
Para garantizar la precisión terminológica de esta investigación, se definen a continuación los conceptos centrales que estructuran el problema, organizados desde el fenómeno tecnológico general hasta las variables específicas de medición.

\begin{itemize}
    \item \textbf{Billetera digital y aplicación financiera móvil:} Herramienta de software instalada en un dispositivo móvil que permite a los usuarios almacenar dinero electrónicamente, realizar transferencias de fondos y ejecutar pagos digitales sin necesidad de intermediación física o dinero en efectivo \cite{bpmconsulting}. En el contexto de este estudio, este término se limita estrictamente a las plataformas Nequi, DaviPlata y la aplicación móvil de Bancolombia, excluyendo el análisis de pasarelas de pago web, criptomonedas o aplicaciones de entidades bancarias internacionales.
    
    \item \textbf{Dependencia funcional tecnológica:} Grado en el que un usuario necesita de manera indispensable un sistema de información para ejecutar tareas regulares, de modo que la ausencia del sistema imposibilita la acción \cite{saviasalud}. Para este proyecto, se define como la proporción de pagos diarios que el estudiante realiza exclusivamente a través de Nequi, DaviPlata o Bancolombia, y la incapacidad percibida para sustituir este medio de pago en su rutina universitaria inmediata.
    
    \item \textbf{Interrupción operativa del servicio (Caída o falla):} Evento en el cual la infraestructura tecnológica de la plataforma financiera (servidores, bases de datos o canales de comunicación) presenta latencias prolongadas, congelamiento de transacciones o inaccesibilidad total \cite{cyberhoot}. Para los fines de esta investigación, se consideran únicamente las fallas atribuibles a la entidad financiera, excluyendo interrupciones causadas por falta de datos móviles, descarga de batería o pérdida de contraseñas por parte del estudiante.
    
    \item \textbf{Actividades cotidianas (Contexto universitario):} Conjunto de transacciones de bajo monto y alta frecuencia necesarias para el desarrollo de la jornada presencial del estudiante. En este estudio, se delimitan específicamente a tres categorías dentro y en los alrededores de la UPTC Sede Tunja: 1) Alimentación (compra en cafeterías y restaurantes aledaños), 2) Apoyo académico (pago de fotocopias, impresiones y materiales), y 3) Movilidad (pago de transporte público urbano o colectivo). Se excluyen transacciones mayores como el pago de matrículas universitarias.
    
    \item \textbf{Mecanismo de contingencia:} Estrategia de respaldo o ruta alternativa que adopta un usuario u organización para mantener su operatividad frente a un evento de falla \cite{ibm_contingencia}. En este estudio, refiere a las acciones inmediatas que toma el estudiante ante la caída de su aplicación preferida, limitadas a: uso de efectivo de reserva guardado previamente, migración temporal a una billetera digital secundaria, préstamo de dinero por parte de un par, o la cancelación definitiva de la compra.
\end{itemize}
\begin{center}
    {\Large\bfseries 8. DISEÑO METODOLÓGICO}
\end{center}
\enconstruccion
\begin{center}
    {\Large\bfseries 9. RESULTADOS Y DISCUSIÓN}
\end{center}
\enconstruccion

\begin{center}
    {\Large\bfseries 10. CONCLUSIONES Y RECOMENDACIONES}
\end{center}
\enconstruccion


\bibliographystyle{ieeetr}
\bibliography{referencias}

% ==========================================
% ANEXOS
% ==========================================
\begin{center}
    {\Large\bfseries ANEXOS}
\end{center}
\enconstruccion

\end{document}